%% kapa: exam=1 round=31 year=2020 period=1 form=A subject=economics subject_ko=경제학원론 seq=39 number=79
%% source: Q-Net 감정평가사 기출문제 (cst00309) artlSeq=5207219 file=extracted/1008/5207219/A형/A형/2020년도 제31회 감정평가사 1차 시험_1교시_A형.pdf
%% conversion: kordoc 4.24.51 -> tools/gen1.py (auto); review=proofread-sonnet; flags=none
\begin{question}{79}
폐쇄경제 IS-LM 모형에서 물가 수준이 하락할 경우 새로운 균형에 관한 설명으로 옳은 것을 모두 고른 것은? (단, 초기 경제는 균형 상태이며, IS 곡선은 우하향, LM 곡선은 우상향)
\begin{framedtext}
ㄱ. 명목 이자율이 하락한다. \par
ㄴ. 투자가 감소한다. \par
ㄷ. 명목 통화량이 증가한다. \par
\end{framedtext}
\begin{choices}
  \item ㄱ
  \item ㄴ
  \item ㄱ, ㄷ
  \item ㄴ, ㄷ
  \item ㄱ, ㄴ, ㄷ
\end{choices}
\end{question}
\answer{1}
